\section{引言}\label{sec:intro}

农业机器人在田间执行巡检、植保、收获运输和自主导航时，需要同时理解作物、垄沟、土壤、
可通行区域与突发障碍物的空间关系。现有视觉导航系统通常采用模块化流程：
先通过作物行识别、道路或可通行区域分割得到二维结果，再经几何拟合、
坐标变换和规则后处理生成导航约束\cite{nkwochaComprehensiveReviewSensing2025,baiVisionbasedNavigationGuidance2023}。
该流程有两个实际限制:
其一，二维感知、几何恢复和控制坐标变换彼此分离，前级预测误差会逐级传播；
其二，作物品种、田块结构或机器人平台变化后，感知模型和后处理规则往往需要重新设定。
对于具有密集植株、连续垄行和复杂遮挡的农田，仅以二维像素级结果作为导航接口也难以表达高度、
遮挡后空间与株间自由间隙。

三维语义占用预测以机器人自身坐标系中的规则体素网格统一表示环境，
每个体素给出占据状态或语义类别。它能够将自由空间、地面、作物和障碍物以统一接口
提供给后续通行性分析、碰撞检查和视觉导航。
自动驾驶研究已证明，Lift--Splat--Shoot 等显式提升方法\cite{philion2020lift,huang2021bevdet,li2022bevdepth}，
以及 DETR3D、PETR 和 BEVFormer 等查询式方法\cite{wang2022detr3d,li2022petr,li2022bevformer}，
都可从多相机图像学习几何一致的鸟瞰表征；该表征还能扩展至稠密占用\cite{wei2023surroundocc,wang2023openoccupancy}、
时序场景预测和规划\cite{zheng2024occworld,mei2025irwm}。
这些方法常设定六路环视相机、面向全向道路交通场景，并使用较大时空范围与模型规模。
而农业机器人则通常配备低算力平台，多在相对静态、前向作业的田间环境中运行。

因此本文聚焦农业机器人最常见的前向作业视野：输入正前、左前和右前三路相机，
预测前方 $20\,$m 的三维语义占用。然而，农业领域因受到作物生长阶段、地块变化、传感器配置等限制，
难以形成如 NuScenes 这类规模化的自动驾驶数据集。现有研究往往也只针对一种作物，难以进行跨地域、跨作物种类的泛化性研究。
为建立全面、多样、可控的训练与评测环境，
本文首先基于 Unreal Engine 5 构建了 AgriOcc-Dataset：包含 298 个完整序列、38,862 帧多视角样本，
涵盖 10 种作物、3 个生长阶段、5 个时段及两类天气条件的农业语义占用数据。

现有语义占用方法在农业前向视野中仍面临三项关键挑战。
第一，不同 BEV 单元被三路相机直接观测的程度不同；仅依赖局部采样时，
弱可见区域难以获得可靠区域的结构化语义上下文。第二，
近场与远场的物理采样密度和语义细节存在显著差异，
维持全场稠密查询会将大量计算耗费在低细节远场。第三，作物边界、
植株高度变化与小尺度障碍具有强烈的三维局部结构，
将每个 BEV token 独立映射为高度 logits 的统一 MLP 容易平滑这些高梯度体素。
为此，本文提出 AgriOcc，如图\ref{fig:overview}所示。
整体框架由三相机图像编码器、GeoVis-BEV 编码器和 Agri-AMoE 三维占用解码器构成。
图像编码器先从三路图像提取多尺度特征，再通过相机投影将其聚合到机器人前方的 BEV 查询中。
GeoVis-BEV 编码器面向农业前向视野组织 BEV 表征学习：
其中，GVAD 将局部可变形图像采样与可见性加权的全局锚点上下文协同建模，
使弱观测单元能够从可靠可见区域获得结构化补全；
近远场查询布局则保持近场稠密计算，并对远场采用规则稀疏采样与特征恢复，
从而匹配不同距离区域的感知需求与计算预算。
最后，Agri-AMoE 将编码后的 BEV token 提升为三维特征体，并依据平面、
垂直和大范围上下文的结构差异自适应路由专家，端到端输出前方空间的六类语义体素。

本文的贡献如下：
\begin{enumerate}[leftmargin=1.8em]
  \item 构建面向农业机器人三维语义占用的仿真数据集AgriOcc-Dataset，在机器人坐标系下统一预测六类前方环境体素；
  \item 提出 GeoVis-BEV 编码器，将几何可见锚点可变形注意力（GVAD）和近远场查询布局纳入统一的 BEV 表征学习过程：前者以可见锚点上下文补全弱观测区域，后者以近密远疏的查询分配匹配农业前向视野的距离差异；
  \item 提出 Agri-AMoE，将 BEV 表征提升为三维特征体，并以梯度能量自适应路由局部平面、垂直结构和大范围上下文专家；
\end{enumerate}

下文首先回顾相关工作，随后给出任务定义与模型设计，最后报告实验结果，并讨论局限与后续研究方向。
